\begin{lemma}[Симметрическая разность паросочетаний]\label{lem:symdiff}
Пусть \(M_1,M_2\) — паросочетания графа \(G\), \(H=(V,\,M_1\mathbin{\triangle}M_2)\).
Тогда каждая компонента \(H\) — изолированная вершина, простой путь или
простой цикл чётной длины, и вдоль каждого пути и цикла рёбра \(M_1\) и \(M_2\)
чередуются.
\end{lemma}
\begin{proof}
Каждая вершина инцидентна не более чем одному ребру \(M_1\) и не более чем
одному ребру \(M_2\), поэтому \(d_H(v)\le2\), и если \(d_H(v)=2\), то два ребра
при \(v\) из разных паросочетаний. Граф со степенями \(\le2\) — объединение
путей и циклов. Чередование означает, что цикл имеет чётную длину.
\end{proof}
\begin{graphillustration}{Компоненты $M\triangle M'$: чётный цикл и пути. Толстые рёбра — $M$, штриховые — $M'$. В пути справа рёбер $M'$ больше; это $M$-дополняющий путь.}{fig:01-symdiff}
\begin{tikzpicture}[font=\small]
\begin{scope}
\node[gi vertex] (c1) at (0,1) {};
\node[gi vertex] (c2) at (1.2,1) {};
\node[gi vertex] (c3) at (1.8,0) {};
\node[gi vertex] (c4) at (1.2,-1) {};
\node[gi vertex] (c5) at (0,-1) {};
\node[gi vertex] (c6) at (-0.6,0) {};
\draw[gi match] (c1) -- (c2); \draw[gi alternate] (c2) -- (c3);
\draw[gi match] (c3) -- (c4); \draw[gi alternate] (c4) -- (c5);
\draw[gi match] (c5) -- (c6); \draw[gi alternate] (c6) -- (c1);
\node[gi note] at (0.6,-1.7) {чётный цикл: поровну};
\end{scope}
\begin{scope}[shift={(3.6,0)}]
\node[gi vertex] (p1) at (0,0.6) {};
\node[gi vertex] (p2) at (1.1,0.6) {};
\node[gi vertex] (p3) at (2.2,0.6) {};
\draw[gi match] (p1) -- (p2); \draw[gi alternate] (p2) -- (p3);
\node[gi note] at (1.1,-0.3) {чётный путь:\\поровну};
\end{scope}
\begin{scope}[shift={(7.2,0)}]
\node[gi free] (q1) at (0,0.6) {};
\node[gi vertex] (q2) at (1.1,0.6) {};
\node[gi vertex] (q3) at (2.2,0.6) {};
\node[gi free] (q4) at (3.3,0.6) {};
\draw[gi alternate] (q1) -- (q2); \draw[gi match] (q2) -- (q3);
\draw[gi alternate] (q3) -- (q4);
\node[gi note] at (1.65,-0.3) {нечётный путь: $M'$ на 1 больше,\\концы свободны в $M$};
\end{scope}
\end{tikzpicture}
\end{graphillustration}


\begin{theorem}[Берж, 1957]
Паросочетание \(M\) максимально (\(|M|=\alpha'(G)\)) тогда и только тогда,
когда в \(G\) нет \(M\)-дополняющих путей.
\end{theorem}
\begin{proof}
\(\Rightarrow\). Пусть есть \(M\)-дополняющий путь \(P=a_1a_2\ldots a_{2k}\).
Рёбра \(a_2a_3,\ldots,a_{2k-2}a_{2k-1}\) лежат в \(M\), а
\(a_1a_2,a_3a_4,\ldots,a_{2k-1}a_{2k}\) — нет. «Инвертируем» путь:
\[
M'=M\mathbin{\triangle}E(P)=\bigl(M\setminus E(P)\bigr)\cup\bigl(E(P)\setminus M\bigr).
\]
Внутренние вершины пути по-прежнему покрыты ровно одним ребром, концы
\(a_1,a_{2k}\) были свободны и стали покрыты, остальные вершины не затронуты.
Значит, \(M'\) — паросочетание и \(|M'|=|M|+1\). Противоречие с максимальностью.

\(\Leftarrow\). Пусть \(M\)-дополняющих путей нет, а \(M'\) — максимальное
паросочетание. По лемме~\ref{lem:symdiff} граф \(H=(V,M\mathbin{\triangle}M')\) —
объединение чётных циклов и путей с чередующимися рёбрами. В цикле и в пути
чётной длины рёбер \(M\) и \(M'\) поровну. Если бы \(|M'|>|M|\), нашлась бы
компонента-путь, где рёбер \(M'\) больше: путь нечётной длины, начинающийся и
кончающийся рёбрами \(M'\setminus M\). Его концы не покрыты \(M\): конец пути
имеет степень 1 в \(H\), и если бы он был покрыт ребром \(e\in M\), то
\(e\notin M'\) (у вершины уже есть ребро \(M'\)), т.\,е. \(e\in H\) и степень была бы 2.
Получен \(M\)-дополняющий путь — противоречие. Следовательно, \(|M|=|M'|\).
\end{proof}
\begin{graphillustration}{Инверсия дополняющего пути $a_1\ldots a_6$: рёбра пути меняют статус, свободные концы становятся покрытыми, $|M|$ растёт на 1.}{fig:01-augment}
\begin{tikzpicture}[font=\small]
\foreach \i/\x in {1/0,2/1.1,3/2.2,4/3.3,5/4.4,6/5.5} {
  \ifnum\i=1 \node[gi free] (a\i) at (\x,0) {$a_{\i}$};
  \else\ifnum\i=6 \node[gi free] (a\i) at (\x,0) {$a_{\i}$};
  \else \node[gi vertex] (a\i) at (\x,0) {$a_{\i}$}; \fi\fi
}
\draw[gi edge] (a1) -- (a2); \draw[gi match] (a2) -- (a3);
\draw[gi edge] (a3) -- (a4); \draw[gi match] (a4) -- (a5);
\draw[gi edge] (a5) -- (a6);
\node[gi note] at (2.75,-0.8) {$M$: 2 ребра на пути};
\node[gi note] at (6.6,0) {$\Longrightarrow$};
\begin{scope}[shift={(7.6,0)}]
\foreach \i/\x in {1/0,2/1.1,3/2.2,4/3.3,5/4.4,6/5.5} \node[gi vertex] (b\i) at (\x,0) {$a_{\i}$};
\draw[gi match] (b1) -- (b2); \draw[gi edge] (b2) -- (b3);
\draw[gi match] (b3) -- (b4); \draw[gi edge] (b4) -- (b5);
\draw[gi match] (b5) -- (b6);
\node[gi note] at (2.75,-0.8) {$M\triangle E(P)$: 3 ребра};
\end{scope}
\end{tikzpicture}
\end{graphillustration}


\begin{cor}[Алгоритмическая идея]
Начиная с пустого паросочетания и увеличивая его вдоль дополняющих путей,
пока они есть, получаем максимальное паросочетание: процесс заканчивается
не позже чем через \(n/2\) шагов, а по теореме Бержа остановка возможна
только на максимальном паросочетании.
\end{cor}

\begin{example}
На рис.~\ref{fig:notation-matchings} паросочетание \(\{bc\}\) максимально по
включению, но путь \(a,b,c,d\) дополняющий; инверсия даёт \(\{ab,cd\}\).
Именно поэтому жадное добавление рёбер не гарантирует максимальность.
\end{example}
